\subsection*{三条历史从同一个成功开始}

重新使用章首的共同初态，完整执行到 X1、O1、EV1 和 F1=501 都已经提交。现在分别施加三种故障，每一条历史都从这一状态独立出发。

第一条只杀死购买数据库进程，磁盘及恢复所需日志保留，已确认提交满足配置的持久性条件。重启以后，数据库通过恢复记录重建事务一致状态。恢复者沿 S 找到 C1，沿 C1 找到原 O1 和结果，平台余额为 7000。这个故障改变了进程的运行状态，恢复所用的业务历史包含原提交。

第二条提升另一副本。若被提升者完整保存了所有需要保留的已确认提交，并在开放以前把它们恢复为可读状态，购买链可以连续接上。第十一章已经分开说明 WAL 到达、持久刷新和应用的区别。此次切换还要撤销旧主资格，让旧主恢复以后接回新主的历史。

第三条让接班者只保留旧前缀。它可能来自异步副本丢失已响应后缀，也可能来自主动恢复一天以前的备份。当前购买库中 S、C1、O1、EV1 都可能消失，E1 的登记和下一编号也回到过去。平台 X1 与履约 F1 处于另一个保存域，继续存在。服务若根据新购买库的空查询重新受理，便把已经发生过的购买解释成从未开始。

\CFourteenFigure{ch14-original03}{相同成功初态分别经历进程崩溃、无损接班与旧历史回退，外部 X1 和 F501 在各条历史中独立保存}{fig:ch14-three-recovery}

数据库恢复的内部一致性与业务历史的完整性需要分别判断。一份旧备份可以满足全部表约束，所有外键都有效，事务也没有写到一半，却缺少昨天已经成功的购买。它正确表达了昨天更早时刻的购买库。支付平台继续保存 7000 分余额，消费者继续保存已建的 F501。

PostgreSQL 的连续归档恢复以物理基础备份和连续 WAL 重建数据库状态，可以选择在某个旧时点停止回放~\cite{pg18-recovery}。这给出了旧历史回退的一条真实路径。逻辑备份也可以恢复较早的数据和编号状态，\code{pg\_dump} 产生的逻辑转储与物理基础备份分别使用自己的恢复机制\footnote{官方文档明确说明：“Such dumps are logical and do not contain enough information to be used by WAL replay.” 此处的 logical 指逻辑转储~\cite{pg18-recovery}。}。

\subsection*{旧数据库相同，平台结局可以不同}

为了看清当前读取缺了什么，取一份尚无 S 的备份 B0，构造两条历史。历史 H0 中，S 从未到达，直接恢复 B0，平台余额为 10000。历史 H1 中，S 已经过 A 和 E，平台完成 X1、余额为 7000，再恢复 B0。当前购买库在两条历史中完全相同，都是没有 S，也没有 C1。

对 H0，一次具有有效资格的新受理原本可以开始购买。对 H1，恢复目标是找回原 C1/X1。给两条历史都创建新的平台键，会使 H1 多出一笔付款。给两条历史都返回原成功，又会在 H0 虚构付款。仅依赖 B0 的处理程序接收到相同输入、读到相同状态，无法给出这两种不同结论。

\CFourteenFigure{ch14-original04}{相同旧备份与相同 S 对应未发生付款和已经发生 X1 的两条历史，恢复决定需要备份以外的信息}{fig:ch14-indistinguishable}

增加一个地域代次并没有补回这个区别。H0 和 H1 都可以取得第 9 代。若把重试改成 \code{pay-C9}，H1 就获得另一次平台首次支付机会。第 9 代可以授权新地域接收新工作，可以帮助新编号避开旧编号，却没有描述第 8 代已经发生的 X1。

\subsection*{编号重用怎样把两次购买合成一次}

再看另一种后果。B0 不仅缺少 C1，还保存着下一个支付编号为 1。原历史中，S1 已使用 \code{C1/pay-C1}，平台保存 X1。恢复 B0 以后，Alice 明确发起另一项合法购买 S2，新的本地分配器又给它 C1。若平台输入包含账户 1、C1、CNY 3000，商品规则也相同，两次平台输入可能完全相同。

平台收到新的 \code{pay-C1}，按原规则返回 X1。新库于是可能把 S2 绑定到原来属于 S1 的收款，创建一张表面成功的新购买。Alice 要买两份，账本却只收到一份的钱，原 S1 又失去了本地记录。平台准确合并了相同平台身份，本地错误地把两项业务意图配到了这个身份。

让 S2 的金额改为 3500，平台会在同一 \code{pay-C1} 上发现异参并返回冲突。此时第二项合法购买无法付款，冲突也不能靠覆盖平台原输入解决。若恢复程序为了绕过冲突换成另一个键，它又必须先分清这是 S2 的新购买，还是 S1 的未知重试。

\CFourteenFigure{ch14-original05}{编号计数回退以后，新 S2 重用 C1，相同平台输入误并到 X1，不同金额则在 pay-C1 上冲突}{fig:ch14-number-reuse}

会话登记有同样的历史责任。假如 B0 保存下一会话编号为 1，之后 E1 已签发、关闭并回收，恢复以后再次签发 E1，就会让旧许可和旧 K 进入一份新的活动登记。第十二章依靠缺失旧会话拒绝来接替逐键历史，这次编号重用却重新打开了相同的准入范围。

新分配器因此从未回退的权威来源取得代次。本例目录对所有逻辑分片使用同一份代次分配进度，已经授予任一分片的第 9 代，之后也不再授予其他分片。因此，共用同一商户平台范围的不同分片，各自取得的代次也不同。这个全局来源与本代的持久计数共同承担编号唯一性。

例如第 9 代的新支付编号可以写为 \code{C9:1}，新登记写为 \code{E9:1}。各类编号在本代内单独递增，分配相同类型的多个入口共同仲裁这份计数。固定分隔和编码规则要保证不同字段组合不能拼成相同编号，计数到达上限便停止分配。权威目录持久保留所有已授予的代次，恢复和重新分配都继续使用这份进度。恢复旧购买时继续读原 \code{C1/pay-C1}，新编号规则只作用于真正的新受理。

如果连权威目录也恢复到旧备份，它的第 9 代分配进度同样可能重用。把这份状态放到另一个进程，只改变保存位置。后续协议明确把目录、独立受理日志和平台账本列为本次本地回退以后继续保存的来源。它们自身的复制、备份、灾难恢复和权限管理承担相应的持久性前提。

\subsection*{新时间线与旧外部事实}

PostgreSQL 在归档恢复完成后创建新的时间线，用来区分恢复后生成的 WAL 与原分支，时间线历史帮助恢复者选取相应日志~\cite{pg18-recovery}。例如原时间线已经经过 A、R，新的数据库分支却从 A 以前开始。不同 WAL 分支得以保留，平台 X1 的身份和存在性则由平台账本与业务关联描述。

\CFourteenFigure{ch14-timeline-fork}{数据库从 B0 分出新时间线，平台与履约保存域继续保留原来已经发生的效果}{fig:ch14-timeline-fork}

故障切换文档同时要求旧主恢复时获知自己已经失去主资格，并说明故障判断、通知与完整故障切换系统由部署提供~\cite{pg18-failover}。对本书的支付链，还要再追踪一条越过数据库边界的消息：旧工作者在失去资格以前发出的 \code{pay-C1}，可能已经离开进程，停在网关或平台队列里。结束旧进程与这份请求完成传输是两个事件。

若平台继续保护 \code{pay-C1}，这份迟到到达会回到原 X1 或原未决操作。若平台已遗忘该身份，又允许迟到到达首次执行，本地的新代次阻止不了它。恢复流程需要把旧请求的受理与终结接到平台实际提供的能力，随后才能解释未查到资金记录的含义。
\FloatBarrier
